\section{Fixed arity and pure ownership}
\label{sec:arity}
\label{arity:appendix}

The projective construction of Theorem~\ref{thm:projective} extends to every
fixed arity $r\geq 3$ by
replacing binary triples with binary circuits of size $r$.  This appendix
gives the additional geometry and records the exact finite bounds.  It also
gives a separate realization with one pure owner per resource.  The resource
count is exponential in the length of the construction in both realizations;
the pure realization has exponentially many owners as well.

Throughout, payments are arbitrary nonnegative functions on the entire fixed
real report domain.  For a report profile $b=(b_i)_{i\in N}$, write
$b_C=\sum_{i\in C}b_i$.  The hypotheses used below are coalition coverage and
a uniform funding bound,
\begin{equation}
  \sum_{i\in C}y_i(b)\geq F(b_C)\quad(C\subseteq N),
  \qquad
  \sum_{i\in N}y_i(b)\leq\kappa F(b_N),
  \label{arity:contract}
\end{equation}
at every feasible profile, together with size-$s$ coalition withholding
monotonicity: if $|C|\leq s$, $b'_i\geq b_i$ coordinatewise for $i\in C$,
and $b'_i=b_i$ for $i\notin C$, then
\begin{equation}
  \sum_{i\in C}y_i(b')\geq\sum_{i\in C}y_i(b).
  \label{arity:tcwm}
\end{equation}
Members of $C$ may have unchanged reports.  Coverage in
\eqref{arity:contract} is required for \emph{every} coalition, including
coalitions larger than $s$.  No continuity, formula for payments, anonymity,
or zero payment to a zero reporter is assumed.  Since full OUT implies
\eqref{arity:tcwm}, every lower bound below also applies under full OUT.

\subsection{The binary-circuit construction}
\label{arity:construction-section}

Fix integers $k\geq2$ and $M\geq1$, and put $r=k+1$.  Let
\[
  V=\mathbb F_2^{\,1+(k-1)M}
\]
have the specified basis
\[
  v,\qquad e_{j,\ell}\quad
  (1\leq j\leq M,\ 1\leq\ell\leq k-1).
\]
Resource rows are indexed by $V\setminus\{0\}$, so their number is
\begin{equation}
  d=2^{1+(k-1)M}-1.
  \label{arity:row-count}
\end{equation}
A binary circuit means a nonempty set of distinct rows whose sum over
$\mathbb F_2$ is zero and none of whose proper nonempty subsets has sum zero.
The activities are all binary circuits of cardinality $k+1$.  Each activity
has profit one and consumes one unit of each of its $k+1$ rows.  Thus, with
$\mathcal A$ denoting the set of activities, the full technology is
\begin{equation}
  F(c)=\max\left\{
    \sum_{A\in\mathcal A}x_A:
    x_A\geq0,\quad
    \sum_{\substack{A\in\mathcal A\\u\in A}}x_A\leq c_u
       \text{ for every }u\in V\setminus\{0\}
  \right\}.
  \label{arity:technology}
\end{equation}
Every activity has arity exactly $r$, and every $g\in\operatorname{GL}(V)$
permutes the resources and the activities of this finite technology.

For the first realization take the fixed actual owner set
\[
  N=\{p\}\ \cup\
  \{h_{j,\ell}:1\leq j\leq M,\ 1\leq\ell\leq k-1\}
  \ \cup\ \{\ell_j:1\leq j\leq M\}.
\]
There are $n=1+kM$ owners.  Every owner has the all-ones endowment and the
same fixed report cube $[0,1]^d$.  Let $\delta_u\in\R^d$ denote one unit of
resource $u$, to distinguish real resource quantities from the binary
vectors that index resources.  Set
\[
  w_j=v+\sum_{\ell=1}^{k-1}e_{j,\ell}\quad\text{in }V.
\]
Basis independence makes all displayed root, high and low rows nonzero and
pairwise distinct.
For $1\leq m\leq M$, define a feasible profile $C_m$ by
\begin{equation}
  b_p=\delta_v,\qquad
  b_{h_{j,\ell}}=\delta_{e_{j,\ell}},\qquad
  b_{\ell_j}=
  \begin{cases}
    \delta_{w_j},&j\leq m,\\
    0,&j>m.
  \end{cases}
  \label{arity:profiles}
\end{equation}
All actual owners and all their report domains remain fixed as $m$ changes.
In particular, the high owners of inactive branches already report their
basis rows.

\begin{lemma}[Supported activities]
\label{arity:supported}
At $C_m$, the only activities whose rows all have positive capacity are
\[
  A_j=\{v,e_{j,1},\ldots,e_{j,k-1},w_j\},
  \qquad 1\leq j\leq m.
\]
Consequently $F(b_N(C_m))=1$.
\end{lemma}

\begin{proof}
Consider a nonempty zero-sum subset of the larger set
\[
  S=\{v\}\cup\{e_{j,\ell}:j,\ell\}\cup\{w_j:1\leq j\leq M\}.
\]
Let $J$ be its set of selected low-row indices and put $t=|J|$.  The basis
coordinate $e_{j,\ell}$ occurs in only two members of $S$: the row
$e_{j,\ell}$ itself and the row $w_j$.  Cancellation therefore requires all
$k-1$ high rows of every branch in $J$, and no high rows of any other branch.
Cancellation of the $v$ coordinate requires the root row precisely when $t$
is odd.  If $t=0$, these conditions give the empty subset.  Every nonempty
zero-sum subset thus has cardinality
\begin{equation}
  kt+(t\bmod 2).
  \label{arity:dependency-cardinality}
\end{equation}
For $k\geq2$, this cardinality equals $k+1$ precisely when $t=1$: if $t\geq2$
it is at least $2k>k+1$.  The corresponding set is $A_j$, and it is minimal
because its $k$ rows $v,e_{j,1},\ldots,e_{j,k-1}$ are independent and $w_j$
is their sum.  At $C_m$ its low row has positive capacity exactly when
$j\leq m$, which proves the classification.

Running any one active branch at level one is feasible.  Every supported
activity consumes the root row, whose capacity is one, so production cannot
exceed one.  Equivalently, a dual certificate for the \emph{full} technology
assigns price one to $v$ and to every zero-capacity row, and price zero to
the other rows.  A supported activity contains $v$; every other activity
contains a zero-capacity row.  Each dual constraint is therefore satisfied,
and the dual objective is one.
\end{proof}

\begin{lemma}[Seed and private-owner orbits]
\label{arity:orbits}
After lawful finite averaging over resource automorphisms and owner
permutations, the owners of the first branch at $C_1$ have a common payment
$q_1\geq1/(k+1)$.  At $C_m$, all $km$ active private owners
\[
  \{h_{j,\ell}:j\leq m,\ 1\leq\ell\leq k-1\}
  \cup\{\ell_j:j\leq m\}
\]
have a common payment $q_m$.
\end{lemma}

\begin{proof}
The finite averaging argument of Lemma~\ref{lem:average} applies to
$\operatorname{GL}(V)\times\mathfrak S_N$.  A binary linear automorphism
acts on a real report by \emph{permuting its resource coordinates}; it does
not take binary linear combinations of real capacities.  These resource
permutations preserve \eqref{arity:technology}, the real report cubes and
coordinatewise order.  Owner permutations preserve the common fixed report
domains, coalition sizes and the coverage and comparison inequalities.
Consequently conjugating a rule and averaging preserves nonnegativity,
\eqref{arity:contract}, \eqref{arity:tcwm}, and its funding factor.  It
suffices to prove the lower bound for this equivariant averaged rule.

At $C_1$, any permutation of the $k$ basis vectors
$v,e_{1,1},\ldots,e_{1,k-1}$ fixes their sum $w_1$ and fixes every high
basis row of an inactive branch.  Also the map
\[
  v\longmapsto w_1,\qquad e_{j,\ell}\longmapsto e_{j,\ell}
       \quad\text{for all }j,\ell
\]
is invertible: the coefficient of $v$ in the replacement for $v$ is one.
It exchanges $v$ and $w_1$.  These maps preserve the entire multiset of
reported rows at $C_1$, including the inactive high reports and zero low
reports.  Combined with the matching owner permutations, they stabilize
$C_1$ and act transitively on its $k+1$ first-branch owners.  Their coalition
can run $A_1$ at level one, so coverage gives $(k+1)q_1\geq1$.  This use of
coverage does not require $k+1\leq s$.

At $C_m$, permute the active branch blocks, or permute the high basis vectors
within one active branch, fixing $v$ and the inactive basis vectors.
Additionally, for each $j\leq m$ the basis replacement
\[
  e_{j,1}\longmapsto w_j
\]
with every other basis vector fixed is invertible, since the coefficient of
$e_{j,1}$ in $w_j$ is one.  It fixes $v$, exchanges $e_{j,1}$ with $w_j$,
and fixes all other displayed rows except these two.  Within-branch
permutations and this exchange connect all $k$ private rows of a branch;
active branch permutations connect the branches.  They preserve the entire
profile, so the induced stabilizers are transitive on all $km$ active
private owners.  Equivariance gives their common payment $q_m$.
\end{proof}

\begin{theorem}[Finite fixed-arity lower bound]
\label{arity:finite-bound}
Fix integers $k\geq2$, $s\geq2$ and $M\geq1$.  There is a finite institution
of arity exactly $k+1$, with $1+kM$ fixed fully eligible owners and
$2^{1+(k-1)M}-1$ resource rows, for which every payment rule satisfying
\eqref{arity:contract}--\eqref{arity:tcwm} has funding factor at least
\begin{equation}
  L_{k,M,s}
  :=1+\frac{k(M-1)}{k+1}P_{k,M,s},\qquad
  P_{k,M,s}:=\prod_{m=2}^{M}\frac{a_m}{a_m+1},\qquad
  a_m:=\min\{s-1,k(m-1)\}.
  \label{arity:finite-tail}
\end{equation}
Empty products are one.  If $s\geq k(M-1)+1$, including $s=1,M=1$, the
same construction gives
\begin{equation}
  \kappa\geq\frac{k}{k+1}Q_{k,M},\qquad
  \kappa\geq1+\frac{k(M-1)}{(k+1)M}Q_{k,M},\qquad
  Q_{k,M}:=\prod_{m=2}^{M}\frac{km}{k(m-1)+1}.
  \label{arity:untruncated}
\end{equation}
\end{theorem}

\begin{proof}
Use \eqref{arity:profiles} and the averaged rule from
Lemma~\ref{arity:orbits}.  In the transition from $C_{m-1}$ to $C_m$, only
the same actual owner $\ell_m$ changes its report, from $0$ to $\delta_{w_m}$.
Select any $a_m$ of the $k(m-1)$ old active private owners and protect them
together with $\ell_m$.  This fixed coalition has $a_m+1\leq s$ members.
All selected old members keep their reports unchanged.  At the old profile
their total payment is $a_mq_{m-1}$, and the restoring owner's old payment
is nonnegative.  At the new profile all selected members, including the
restorer, receive $q_m$.  Hence \eqref{arity:tcwm} implies the exact recurrence
\begin{equation}
  (a_m+1)q_m\geq a_mq_{m-1},\qquad 2\leq m\leq M.
  \label{arity:recurrence}
\end{equation}
The selected old owners need not comprise whole branches.  Iterating
\eqref{arity:recurrence} from $(k+1)q_1\geq1$ gives
\begin{equation}
  q_M\geq\frac{P_{k,M,s}}{k+1}.
  \label{arity:terminal-q}
\end{equation}

Write $z$ for the root owner's payment at $C_M$.  Coverage for the actual
owners of any complete branch and the budget at $C_M$ give respectively
\[
  z+kq_M\geq1,\qquad
  \kappa\geq z+kMq_M.
\]
The second inequality uses Lemma~\ref{arity:supported} and discards every
other owner's payment only by nonnegativity.  Combining them yields
$\kappa\geq1+k(M-1)q_M$, proving \eqref{arity:finite-tail}.  This includes
$M=1$, when terminal coverage alone gives $\kappa\geq1$.

If $s\geq k(M-1)+1$, protect \emph{all} old active private owners at every
transition.  The recurrence becomes
\[
  [k(m-1)+1]q_m\geq k(m-1)q_{m-1}.
\]
Writing $T_m=kmq_m$ gives
\[
  T_1\geq\frac{k}{k+1},\qquad
  T_m\geq\frac{km}{k(m-1)+1}T_{m-1}.
\]
Thus $T_M\geq kQ_{k,M}/(k+1)$, and the budget proves the first bound in
\eqref{arity:untruncated}.  Since
$q_M\geq Q_{k,M}/[(k+1)M]$, terminal branch coverage proves the second.
When $M=1$ there are no comparisons, so this argument also applies at $s=1$.

All comparisons just used occur inside the one original fixed institution.
Symmetry is a proof reduction, not a restriction on admissible payment
rules.  If the original rule is asymmetric, an expensive profile may be an
image of $C_M$ under the finite symmetry group, rather than $C_M$ itself.
The same lower bound on its worst-case funding factor follows.
\end{proof}

The finite bound also has a direct nonnegative rational certificate.
Put $c=k(M-1)$ and $P=P_{k,M,s}$.  Give the terminal coverage inequality
$z+kq_M\geq1$ and the budget inequality
$\kappa-z-kMq_M\geq0$ weight one each.  Give the seed inequality
$(k+1)q_1\geq1$ weight $cP/(k+1)$, and give recurrence $m$ weight
\begin{equation}
  \lambda_m=\frac{c}{a_m+1}
    \prod_{j=m+1}^{M}\frac{a_j}{a_j+1}.
  \label{arity:certificate-weights}
\end{equation}
For $M\geq2$, the coefficient of $q_M$ from the final two inequalities is
$-c$, canceled by $\lambda_M(a_M+1)$.  For $2\leq m<M$,
$\lambda_m(a_m+1)=\lambda_{m+1}a_{m+1}$, and
$\lambda_2a_2=cP$ cancels the seed term.  The $z$ coefficients also cancel,
leaving exactly $\kappa\geq1+cP/(k+1)$.  At $M=1$, $c=0$ and terminal
coverage plus budget already give the conclusion.  These are consequences
for arbitrary payment variables; no marginal-payment representation is used.

\subsection{Length optimization and the fixed-arity exponent}
\label{arity:rate-section}

\begin{corollary}[Optimized length within the construction]
\label{arity:optimized}
For fixed integers $k\geq2$ and $s\geq2$, set
\[
  h=\left\lceil\frac{s}{k}\right\rceil,
  \qquad \rho=\frac{s-1}{s}.
\]
The largest value of $L_{k,M,s}$ over positive integer lengths $M$ is
attained exactly at $M=s$ and $M=s+1$, with common value
\begin{equation}
  B_{k,s}
  =1+\frac{k(s-1)}{(k+1)h}Q_{k,h}\rho^{\,s-h}.
  \label{arity:optimized-bound}
\end{equation}
The length $M=s$ realizes this bound with
$n=1+ks$ fully eligible owners and $d=2^{1+(k-1)s}-1$ rows.  For fixed $k$,
\begin{equation}
  B_{k,s}-1\sim
  \frac{k}{k+1}\Gamma\!\left(1+\frac1k\right)
       k^{1/k}\exp\!\left(-\frac{k-1}{k}\right)
       s^{(k-1)/k}
       \qquad(s\longrightarrow\infty).
  \label{arity:asymptotic}
\end{equation}
This optimizes the displayed certificate over its length; it does not
identify the optimal funding factor among all mechanisms.
\end{corollary}

\begin{proof}
The cutoff satisfies $k(h-1)\leq s-1$ and $kh\geq s$.  Through $m=h$ the
recurrence is untruncated, and for $m\geq h+1$ its factor is $\rho$.
Telescoping the factors $(m-1)/m$ therefore gives
\[
  P_{k,h,s}=\frac{Q_{k,h}}{h},\qquad
  P_{k,M,s}=\frac{Q_{k,h}}{h}\rho^{M-h}\quad(M\geq h).
\]
Let $I_M=(M-1)P_{k,M,s}$.  We have $I_1=0<I_2$.  If $1<M<h$,
\[
  \frac{I_{M+1}}{I_M}
  =\frac{kM^2}{(M-1)(kM+1)}>1,
\]
because the numerator minus the denominator is $(k-1)M+1>0$.  If
$M\geq h$ and $M>1$, the ratio is
\[
  \frac{I_{M+1}}{I_M}=\frac{M\rho}{M-1},
\]
which exceeds one for $M<s$, equals one for $M=s$, and is less than one
for $M>s$.  As $h\leq s$, this proves the two asserted maximizers, including
the boundary case $h=1$.  Substituting $M=s$ proves
\eqref{arity:optimized-bound}.

The gamma functional equation gives the exact identity
\[
  Q_{k,h}
  =\frac{\Gamma(h+1)\Gamma(1+1/k)}{\Gamma(h+1/k)}.
\]
Stirling's formula yields
$Q_{k,h}\sim\Gamma(1+1/k)h^{(k-1)/k}$.  Moreover $h\sim s/k$ and
$\rho^{s-h}\longrightarrow\exp(-(k-1)/k)$.
Inserting these limits into \eqref{arity:optimized-bound} proves
\eqref{arity:asymptotic}.
\end{proof}

\begin{corollary}[An all-$s$ bound]
\label{arity:all-s}
For every fixed integer $r\geq3$ and every integer $s\geq1$, writing
$k=r-1$, the unrestricted worst-case optimal funding factor satisfies
\begin{equation}
  K_r(s)\geq
  \frac{k}{k+1}\left(\frac{s}{2k}\right)^{(k-1)/k}.
  \label{arity:elementary-rate}
\end{equation}
In particular,
$K_r(s)=\Omega_r(s^{(r-2)/(r-1)})$.
\end{corollary}

\begin{proof}
Put $\beta=(k-1)/k$.  For $m\geq2$, concavity of $x^\beta$ gives
\[
  \frac{km}{k(m-1)+1}
    =\frac{m}{m-\beta}
    \geq1+\frac{\beta}{m}
    \geq\left(1+\frac1m\right)^\beta.
\]
Consequently $Q_{k,M}\geq((M+1)/2)^\beta$, also at $M=1$.
Choose $M=\lceil s/k\rceil$.  Then
$k(M-1)+1\leq s$ and $M\geq s/k$, so the first bound in
\eqref{arity:untruncated} implies \eqref{arity:elementary-rate}.
The finite institution is admissible for the worst-case class defining
$K_r(s)$.
\end{proof}

\subsection{A separate realization with pure owners}
\label{arity:pure-section}

\begin{corollary}[Fixed pure-owner realization]
\label{arity:pure}
All the lower bounds above also hold for institutions with one fixed pure
unit owner per resource row.  For the length-$M$ instance, this realization
has
\begin{equation}
  n=d=2^{1+(k-1)M}-1,
  \label{arity:pure-counts}
\end{equation}
and every owner's fixed report domain is the full real interval along its
own resource coordinate.  It does not have the fully eligible realization's
owner count $1+kM$.
\end{corollary}

\begin{proof}
Keep the technology \eqref{arity:technology}.  For each row
$u\in V\setminus\{0\}$ introduce the actual owner $i_u$ with fixed
endowment $\delta_u$ and report domain
\[
  \{t_u\delta_u:0\leq t_u\leq1\}.
\]
All these owners exist at every profile.  A resource automorphism $g$ acts
\emph{jointly} on resource rows and owner labels, sending owner $i_u$ to
owner $i_{g(u)}$.  On scalar reports the action is
\[
  (g\cdot t)_{g(u)}=t_u.
\]
This preserves the fixed institution: the transformed owner still reports
only its own resource.  The associated owner bijection takes a coalition
to another coalition of the same size, maps its total resource report by
the same technology automorphism, preserves coordinatewise restoration,
and leaves all outsiders unchanged in a comparison.  Conjugation therefore
preserves every coverage, funding and incentive inequality.  Averaging over
this finite joint $\operatorname{GL}(V)$ action is lawful.  No independent
owner permutation is required or used.

At the pure-owner version of $C_m$, set $t_u=1$ exactly for
\[
  S_m=\{v\}\cup\{e_{j,\ell}:j,\ell\}
                 \cup\{w_j:j\leq m\},
\]
and set $t_u=0$ for every other row.  Its aggregate capacities coincide with
those of \eqref{arity:profiles}, so Lemma~\ref{arity:supported}, including
the full-technology dual, applies without change.  Each seed or private
stabilizer described in Lemma~\ref{arity:orbits} preserves $S_m$ as a set.
Being a permutation of all nonzero rows, it also preserves its complement.
Its joint action thus fixes the entire report profile, including every
additional zero-report owner.  The same seed and private orbits give
$(k+1)q_1\geq1$ and a common payment $q_m$ for the $km$ active private owners.

In the transition from $C_{m-1}$ to $C_m$, only the same actual pure owner
$i_{w_m}$ restores its report from zero to one.  Protect this owner together
with the specified old active pure owners.  Their labels, eligibility and
reports remain fixed throughout the comparison.  Equation
\eqref{arity:recurrence} follows exactly as before, by discarding only the
restorer's nonnegative old payment.  Terminal branch coverage involves
the $k+1$ actual pure owners of that branch.  All omitted owners, including
additional zero reporters, have nonnegative payments, so the terminal
budget argument is unchanged.  This proves the finite bounds and all their
stated consequences with the counts \eqref{arity:pure-counts}.
\end{proof}

For example, $k=2$ specializes to the arity-three construction.  Its optimized
length-$s$ bound has coefficient
$\sqrt{2\pi}/(3\sqrt e)$ in front of $\sqrt s$, and $B_{2,6}=493/243$.
The length-six witness uses $127$ resource rows in either realization,
with $13$ fully eligible owners or $127$ pure owners.  Neither the
fixed-arity extension nor the change of ownership realization gives a
bound with a fixed resource dimension.
